\documentclass[12pt]{article}

\usepackage[utf8]{inputenc}
\usepackage[romanian]{babel}
\usepackage{amsmath, amssymb, amsthm}
\usepackage{geometry}
\geometry{a4paper, margin=2.5cm}

\title{Cazul Particular în Demersul Didactic}
\author{}
\date{}

\theoremstyle{definition}
\newtheorem{definition}{Definiție}

\theoremstyle{plain}
\newtheorem{proposition}{Propoziție}

\begin{document}

\maketitle

\begin{definition}[Cazul particular]
Cazul particular al unei afirmații matematice sau al unei probleme se obține prin:
\begin{enumerate}
    \item înlocuirea unei variabile cu o constantă;
    \item adăugarea unei condiții la ipoteză;
    \item înlăturarea unei condiții din concluzie.
\end{enumerate}
\end{definition}

\begin{proposition}Roluri ale cazului particular în demersurile didactice și cognitive:
\begin{enumerate}
    \item Pentru elevii de vârstă mică (până la 10–11 ani), intuiția și exemplele variate constituie fundamentul achiziției de cunoștințe și deprinderi matematice.
    \item În general, este recomandat ca prezentarea diverselor tematici din matematică să fie însoțită de un număr semnificativ de exemple relevante (cazuri particulare).
    \item Cazul particular întărește încrederea în enunț.
    \item Cazul particular poate înlătura anumite dubii privitoare la enunț.
    \item Cazul particular poate furniza răspunsul la întrebarea problemei.
    \item Cazul particular poate constitui o etapă a demonstrației (în cadrul așa-numitei demonstrații treptate).
    \item Demonstrația cazului particular poate sugera maniera de abordare.
    \item Cazul particular rezolvă probleme de tipul "există... - afirmație".
\end{enumerate}
\end{proposition}

\end{document}
